% Vietnamese translation for OI Public Library
% Source-PDF: IOI中国国家候选队论文/2009/刘聪/浅谈数位类统计问题.pdf
\documentclass[11pt]{article}
\usepackage{oipl-vn}

\newcommand{\Oh}{\mathcal{O}}
\newcommand{\grlex}{\mathrel{<_{\mathrm{grlex}}}}

\title{Bàn sơ lược về các bài toán thống kê theo chữ số}
\author{Lưu Thông\\Trường Trung học số 2 Thanh Đảo, Sơn Đông}
\date{2009}

\begin{document}
\maketitle

\origtitle{On digit-statistics problems}
\sourcepath{IOI-China-Candidate-Team-Papers/2009/Liu-Cong/digit-statistics.pdf}

\begin{abstract}
Trong thi tin học có một lớp bài toán thống kê trên đoạn liên quan đến chữ số.
Những bài này thường mang màu sắc toán học khá rõ, không thể giải bằng vét cạn
mà cần thực hiện truy hồi và các thao tác tương tự trên từng chữ số. Qua một
vài ví dụ, bài viết giới thiệu ngắn gọn các tư tưởng và phương pháp cơ bản để
giải lớp bài toán này.
\end{abstract}

\noindent\textbf{Từ khóa:} chữ số, đoạn, thống kê, truy hồi, cây, nhị phân.

\section{Mở đầu}

Trong thi tin học có một lớp bài toán như sau: cho một đoạn, yêu cầu tìm một
số ở hệ cơ số $D$ thỏa mãn điều kiện cho trước, hoặc đếm số lượng các số như
vậy. Điều kiện cần xét thường liên quan đến chữ số, chẳng hạn tổng chữ số, số
lượng một chữ số chỉ định, cách chia nhóm theo thứ tự lớn nhỏ của số, v.v. Đoạn
được cho trong đề thường rất lớn, nên không thể dùng phương pháp thô sơ.

Khi đó ta cần lợi dụng tính chất của chữ số để thiết kế thuật toán có độ phức
tạp cỡ $\log n$. Tư tưởng cơ bản nhất để giải loại bài toán này là phương pháp
\term{xác định từng chữ số}. Dưới đây ta xét một số bài mẫu để hiểu cụ thể hơn
về dạng bài và cách suy nghĩ này.

\section{Ví dụ 1: Amount of degrees (URAL 1057)}

\subsection{Đề bài}

Cho đoạn $[X,Y]$. Hãy đếm số nguyên trong đoạn này thỏa mãn điều kiện: số đó
bằng tổng của đúng $K$ lũy thừa nguyên khác nhau của $B$. Chẳng hạn, với
$X=15$, $Y=20$, $K=2$, $B=2$, chỉ có ba số sau thỏa mãn:
\[
  17=2^4+2^0,\qquad
  18=2^4+2^1,\qquad
  20=2^4+2^2.
\]

Dữ liệu vào gồm hai số nguyên $X,Y$ ở dòng đầu, sau đó là $K$ và $B$. Dữ liệu
ra là một số nguyên, tức số lượng số thỏa mãn. Ràng buộc:
\[
  1\le X\le Y\le 2^{31}-1,\qquad
  1\le K\le 20,\qquad 2\le B\le 10.
\]

\subsection{Phân tích}

Số cần tìm là tổng của các lũy thừa khác nhau, nghĩa là trong biểu diễn cơ số
$B$ của nó, mọi chữ số chỉ có thể là $0$ hoặc $1$. Vì vậy, trước hết ta chỉ
cần thảo luận trường hợp nhị phân; các cơ số khác đều có thể chuyển về một bài
toán nhị phân tương ứng.

Rõ ràng miền dữ liệu rất lớn, không thể liệt kê. Độ phức tạp phải ở cấp
$\log n$, nên ta phải xử lý trên chữ số.

Bài này thỏa mãn tính chất trừ theo đoạn. Gọi $\operatorname{count}[i..j]$ là
số lượng số hợp lệ trong đoạn $[i,j]$, khi đó
\[
  \operatorname{count}[i..j]
  =\operatorname{count}[0..j]-\operatorname{count}[0..i-1].
\]
Nói cách khác, với một $n$ cho trước, ta chỉ cần tính có bao nhiêu số từ $0$
đến $n$ thỏa mãn điều kiện.

Giả sử $n=13$, biểu diễn nhị phân là $1101$, và $K=3$. Mục tiêu là đếm các số
từ $0$ đến $13$ có biểu diễn nhị phân chứa đúng $3$ bit $1$. Để dễ suy nghĩ,
hãy tưởng tượng một cây nhị phân đầy đủ cao $4$, trong đó mỗi cạnh hoặc mỗi
bước đi tương ứng với chọn chữ số $0$ hoặc $1$. Mỗi lá biểu diễn một số nhị
phân $4$ chữ số, tức các số từ $0$ đến $2^4-1$.

Đường đi đến lá biểu diễn $n$ chia những số nhỏ hơn $n$ thành một số cây con
đầy đủ. Mỗi khi trên đường đi ta rẽ phải, toàn bộ cây con bên trái tại vị trí
đó đều gồm các số nhỏ hơn $n$. Do đó để đếm các số nhỏ hơn $13$, ta chỉ cần
đếm trong các cây con đầy đủ này: có cây cần đếm số có $3$ bit $1$, có cây chỉ
cần thêm $2$ bit $1$ vì trên tiền tố đã có một bit $1$, và có cây chỉ cần thêm
$1$ bit $1$. Các cây con có cùng chiều cao thì kết quả thống kê như nhau.
Đừng quên xét chính $n$; trong cài đặt, có thể tăng $n$ lên $1$ ngay từ đầu để
tránh bàn riêng trường hợp này, nhưng phải chú ý tràn số.

Còn lại là bài toán: trong một cây nhị phân đầy đủ cao $i$, có bao nhiêu lá mà
biểu diễn nhị phân chứa đúng $j$ bit $1$? Điều này được tính dễ dàng bằng truy
hồi. Đặt $f[i,j]$ là đáp án. Tách theo cây con trái và phải, ta có
\[
  f[i,j]=f[i-1,j]+f[i-1,j-1].
\]

Vậy thuật toán trả lời là: tiền xử lý $f$, rồi với $n$ đầu vào, đi từ gốc đến
lá tương ứng trong cây đầy đủ tưởng tượng; mỗi khi rẽ phải thì cộng số lượng
số trong cây con trái. Mã C++ trong bài gốc như sau:

\begin{verbatim}
void init() // tien xu ly f
{
    f[0][0]=1;
    for (int i=1;i<=31;++i) {
       f[i][0]=f[i-1][0];
       for (int j=1;j<=i;++j) f[i][j]=f[i-1][j]+f[i-1][j-1];
    }
}

int calc(int x,int k) // dem trong [0..x] cac so co k bit 1
{
     int tot=0,ans=0; // tot: so bit 1 da co tren duong di
     for (int i=31;i>0;--i)
     {
        if (x&(1<<i)) {
             ++tot;
             if (tot>k) break;
             x=x^(1<<i);
        }
        if ((1<<(i-1))<=x) {
             ans+=f[i-1][k-tot];
        }
     }
     if (tot+x==k) ++ans;
     return ans;
}
\end{verbatim}

Vấn đề cuối cùng là xử lý cơ số không phải nhị phân. Với truy vấn $n$, ta cần
tìm số lớn nhất không vượt quá $n$ mà biểu diễn cơ số $B$ chỉ chứa $0$ và $1$:
tìm chữ số đầu tiên từ trái sang phải không phải $0$ hoặc $1$, đổi nó thành
$1$, và đặt tất cả chữ số bên phải bằng $1$. Sau đó xem biểu diễn cơ số $B$ thu
được như một biểu diễn nhị phân rồi trả lời bằng thuật toán trên.

Tiền xử lý truy hồi $f$ có độ phức tạp $\Oh(\log^2 n)$. Có $\Oh(\log n)$ lần
truy vấn, nên tổng độ phức tạp thời gian là $\Oh(\log^2 n)$.

Thực ra mã cuối cùng không thao tác trên cây; ta chỉ dùng hình ảnh cây để tiện
suy nghĩ. Cũng có thể xét trực tiếp từ góc nhìn chữ số: với truy vấn $n$, tìm
một vị trí có chữ số bằng $1$ và đổi nó thành $0$; các chữ số bên phải được
chọn tùy ý, và ta cần đếm trong đó các cách có đúng $K-\mathit{tot}$ bit $1$,
trong đó $\mathit{tot}$ là số bit $1$ ở bên trái vị trí này. Khi đó có thể dùng
công thức tổ hợp. Liệt kê từng vị trí bằng $1$ để cộng dồn là đủ.

Ta thấy $f$ suy ra ở trên chính là các số tổ hợp. Cả hai cách đều dùng tư
tưởng xác định từng chữ số. Khi cảm thấy suy luận trực tiếp trên chữ số hơi
khó, hãy vẽ hình để làm rõ ý tưởng.

\section{Ví dụ 2: Sorted bit sequence (SPOJ 1182)}

\subsection{Đề bài}

Sắp xếp tất cả số nguyên trong đoạn $[m,n]$ theo số lượng bit $1$ trong biểu
diễn nhị phân tăng dần. Nếu hai số có cùng số bit $1$, sắp xếp theo giá trị
của số. Hãy tìm số thứ $k$ trong dãy sau khi sắp xếp. Số âm dùng biểu diễn bù
hai: biểu diễn nhị phân của một số âm cộng với biểu diễn nhị phân của số đối
của nó đúng bằng $2^{32}$.

Ví dụ, với $m=0$, $n=5$, dãy sau khi sắp xếp là:
\[
\begin{array}{c|r|l}
\text{STT} & \text{số thập phân} & \text{biểu diễn nhị phân 32 bit}\\
\hline
1&0&0000\ 0000\ 0000\ 0000\ 0000\ 0000\ 0000\ 0000\\
2&1&0000\ 0000\ 0000\ 0000\ 0000\ 0000\ 0000\ 0001\\
3&2&0000\ 0000\ 0000\ 0000\ 0000\ 0000\ 0000\ 0010\\
4&4&0000\ 0000\ 0000\ 0000\ 0000\ 0000\ 0000\ 0100\\
5&3&0000\ 0000\ 0000\ 0000\ 0000\ 0000\ 0000\ 0011\\
6&5&0000\ 0000\ 0000\ 0000\ 0000\ 0000\ 0000\ 0101
\end{array}
\]

Với $m=-5$, $n=-2$, dãy là:
\[
\begin{array}{c|r|l}
\text{STT} & \text{số thập phân} & \text{biểu diễn nhị phân 32 bit}\\
\hline
1&-4&1111\ 1111\ 1111\ 1111\ 1111\ 1111\ 1111\ 1100\\
2&-5&1111\ 1111\ 1111\ 1111\ 1111\ 1111\ 1111\ 1011\\
3&-3&1111\ 1111\ 1111\ 1111\ 1111\ 1111\ 1111\ 1101\\
4&-2&1111\ 1111\ 1111\ 1111\ 1111\ 1111\ 1111\ 1110
\end{array}
\]

Dữ liệu gồm nhiều bộ test. Dòng đầu là số bộ test, không quá $1000$. Mỗi bộ
test chứa ba số nguyên $m,n,k$. Với mỗi bộ test, in ra số thứ $k$ trong dãy đã
sắp xếp. Ràng buộc:
\[
  mn\ge 0,\qquad
  -2^{31}\le m\le n\le 2^{31}-1,\qquad
  1\le k\le \min\{n-m+1,2147473547\}.
\]

\subsection{Phân tích}

Trước hết xét trường hợp $m,n$ đều dương.

Vì khóa sắp xếp thứ nhất là số bit $1$, khóa thứ hai là giá trị của số, ta dễ
xác định đáp án có bao nhiêu bit $1$: lần lượt thống kê trong đoạn $[m,n]$ số
lượng số có $0,1,2,\ldots$ bit $1$ cho đến khi tổng cộng dồn vượt quá $k$.
Giá trị hiện tại chính là số bit $1$ của đáp án; giả sử là $s$. Thuật toán ở
ví dụ 1 giải được bước thống kê này. Đồng thời, ta cũng biết đáp án là số thứ
mấy trong các số thuộc $[m,n]$ có đúng $s$ bit $1$.

Do đó chỉ cần nhị phân đáp án. Với một giá trị thử $\mathit{ans}$, tính số
lượng số trong $[m,\mathit{ans}]$ có đúng $s$ bit $1$ để quyết định hướng tìm.
Mỗi lần thống kê có độ phức tạp $\Oh(\log n)$, nên phần nhị phân có độ phức
tạp $\Oh(\log^2 n)$; tiền xử lý cũng cùng cấp, vì thế thuật toán khá lý tưởng.

Trường hợp $m<0$ cũng không khó xử lý. Ta có thể bỏ qua bit cao nhất của mọi
số, tìm đáp án rồi đặt bit cao nhất trở lại bằng $1$. Hoặc cũng có thể trực
tiếp xem số âm là số nguyên không dấu $32$ bit và xử lý giống số dương. Cả hai
cách đều cần xử lý riêng trường hợp $n=0$.

\section{Ví dụ 3: Sequence (SPOJ 2319)}

\subsection{Đề bài}

Cho tất cả số nhị phân $K$ bit:
\[
  0,1,\ldots,2^K-1.
\]
Cần chia dãy này thành đúng $M$ nhóm, mỗi nhóm là một đoạn liên tiếp của dãy
ban đầu. Gọi $S_i$ ($1\le i\le M$) là tổng số bit $1$ trong biểu diễn nhị phân
của tất cả số thuộc nhóm thứ $i$, và gọi $S$ là giá trị lớn nhất trong các
$S_i$. Nhiệm vụ là làm cho $S$ nhỏ nhất.

Dữ liệu vào gồm hai số nguyên $K,M$. Dữ liệu ra là giá trị nhỏ nhất của $S$;
đáp án không bảo đảm lưu được trong kiểu số nguyên $64$ bit. Ràng buộc:
\[
  1\le K\le 100,\qquad 1\le M\le 100,\qquad M\le 2^K.
\]

\subsection{Phân tích}

Làm trực tiếp không dễ, nhưng có thể xử lý bằng cách nhị phân đáp án.

Hiển nhiên, nếu biết giá trị $S$, ta có thể tham lam tìm số nhóm ít nhất mà
dãy cần được chia. Khi $S$ tăng, số nhóm này không tăng. Dựa vào tính đơn điệu
đó, ta nhị phân $S$, rồi tham lam tính số đoạn ít nhất cần dùng để phủ
$[0,2^K-1]$ và kiểm tra.

Còn lại là cách tính số nhóm. Vì $M$ nhỏ, mỗi lần ta bắt đầu từ điểm hiện tại
$a$ và tìm $b$ lớn nhất sao cho
\[
  \operatorname{count}(b)-\operatorname{count}(a-1)\le S,
\]
trong đó $\operatorname{count}(i)$ là tổng số bit $1$ trong biểu diễn nhị phân
của mọi số từ $0$ đến $i$. Sau đó lấy $b+1$ làm điểm bắt đầu mới, cho đến khi
một lần tìm được $b$ vượt quá $2^K-1$, hoặc số đoạn đã vượt quá $M$.

Bằng phương pháp xác định từng chữ số, ta dễ tính $\operatorname{count}(i)$.
Tương tự ví dụ 1, trong cây nhị phân đầy đủ cao $K$, đi từ gốc đến lá $i$;
mỗi khi rẽ phải thì cộng trọng số của cây con trái. Nói cách khác, tìm mọi
vị trí bit $1$ của $i$ và cộng trọng số tương ứng; khi cộng phải chú ý số bit
$1$ đã xuất hiện ở tiền tố trước đó.

Cách tìm $b$ cũng không khó. Đặt $\operatorname{find}(i)$ là giá trị $k$ lớn
nhất sao cho tổng số bit $1$ trong các số từ $0$ đến $k$ không vượt quá $i$.
Khi đó
\[
  b=\operatorname{find}(\operatorname{count}(a-1)+S).
\]
Hàm $\operatorname{find}(i)$ cũng dùng phương pháp xác định từng chữ số: đi từ
gốc xuống lá trong cây; nếu trọng số đã có cộng với trọng số cây con trái của
nút hiện tại không vượt quá $i$ thì đi sang phải, ngược lại đi sang trái.

Cả $\operatorname{count}(i)$ và $\operatorname{find}(i)$ đều có độ phức tạp
$\Oh(K)$, nên tổng độ phức tạp thuật toán là $\Oh(MK\log S)$, cộng thêm chi
phí của số học độ chính xác cao.

\section{Ví dụ 4: Tickets (SGU 390)}

\subsection{Đề bài}

Một người bán vé bán vé cho hành khách. Với mỗi hành khách, anh ta bán nhiều
vé liên tiếp, cho đến khi tổng chữ số của các số hiệu vé đã bán không nhỏ hơn
một số dương $k$ cho trước. Sau đó anh ta lặp lại quy tắc đó với hành khách
tiếp theo. Ban đầu người bán vé có các vé mang số nguyên liên tiếp từ $L$ đến
$R$. Hãy tính anh ta có thể bán vé cho bao nhiêu hành khách.

Dữ liệu vào gồm ba số nguyên $L,R,k$. Dữ liệu ra là đáp án. Ràng buộc:
\[
  1\le L\le R\le 10^{18},\qquad 1\le k\le 1000.
\]

\subsection{Phân tích}

Bài này hơi giống ví dụ 3: ta cũng chia các số nguyên trong một đoạn thành các
nhóm liên tiếp. Khác biệt là ở đây $k$ nhỏ, tức số nhóm có thể rất lớn, nên
không thể tính từng nhóm một.

Như trước, ta bắt đầu từ trường hợp đặc biệt. Giả sử đoạn đã cho là
$[1,10^h-1]$, tức một cây thập phân đầy đủ cao $h$. Hãy xét cách ghép các bài
toán con thành bài toán gốc. Khó khăn gặp phải là việc ``nối'' giữa hai cây
con: nhóm cuối của cây con trước chưa chắc đã đầy, nên vài phần tử đầu của cây
con sau có thể bị ghép vào nhóm cuối ấy. Để giải quyết, ta thêm một chiều để
ghi trước mỗi cây con còn bao nhiêu ``chỗ trống''. Từ đó có truy hồi sau.

Đặt $f[h,\mathit{sum},\mathit{rem}]$ là kết quả đối với một cây thập phân đầy
đủ cao $h$, trong đó mỗi số thuộc đoạn này cần được cộng thêm tổng chữ số
$\mathit{sum}$ ở tiền tố, và trước đoạn này nhóm hiện tại còn
$\mathit{rem}$ đơn vị trống. Giá trị trả về của $f$ cần ghi hai số:
$f[h,\mathit{sum},\mathit{rem}].a$ là số nhóm tạo ra, còn
$f[h,\mathit{sum},\mathit{rem}].b$ là phần trống còn lại của nhóm cuối, để có
thể ghép với cây con kế tiếp.

Giá trị $f$ được suy ra từ mười cây con:
\[
\begin{aligned}
\text{for }i:=0\text{ to }9\text{ do }\quad
f[h,\mathit{sum},\mathit{rem}]
 :=\operatorname{merge}\bigl(&f[h,\mathit{sum},\mathit{rem}],\\
&f[h-1,\mathit{sum}+i,
  f[h,\mathit{sum},\mathit{rem}].b]\bigr).
\end{aligned}
\]
Hàm $\operatorname{merge}$ biểu thị ghép hai bản ghi, với giả mã:
\begin{verbatim}
function merge(x,y);
  x.a:=x.a+y.a;
  x.b:=y.b;
  return x;
\end{verbatim}

Khi $h=0$, nếu $\mathit{rem}>0$ thì $f[h,\mathit{sum},\mathit{rem}].a=0$,
ngược lại $f[h,\mathit{sum},\mathit{rem}].a=1$; giá trị $b$ cần tính theo quan
hệ giữa $\mathit{sum}$ và $k$.

Như vậy ta đã tính được số nhóm của các cây con đầy đủ. Công việc tiếp theo
dễ hơn: trong cây thập phân đầy đủ, đi từ lá $L$ đến lá $R$, tính và ghép các
cây con đầy đủ gặp trên đường. Cụ thể, đi từ lá chứa $L$ lên đến tầng ngay
dưới tổ tiên chung gần nhất của $L$ và $R$, trên đường đó tính các cây anh em
ở bên phải. Sau đó, tại tầng ngay dưới LCA, đi từ tổ tiên của $L$ sang tổ tiên
của $R$ và tính các cây con đi qua. Cuối cùng đi xuống lá $R$, trên đường tính
các cây anh em ở bên trái. Tất nhiên cũng không được quên tính chính các lá
chứa $L$ và $R$.

Trong hình minh họa của bài gốc, tác giả dùng cây nhị phân để dễ vẽ, nhưng ý
tưởng hoàn toàn giống cây thập phân: trước hết tìm LCA của $L$ và $R$, đi
ngược từ $L$ lên và cộng các cây con bên phải, rồi đi sang phía tổ tiên của
$R$, cuối cùng đi xuống $R$ và cộng các cây con bên trái.

Cần xử lý riêng các trường hợp $L=R$ và $R=10^{18}$ nếu chỉ tính đến $18$
tầng. Ngoài ra, nếu nhóm cuối cùng chưa đầy thì không tính vào tổng số nhóm;
điểm này cần đặc biệt chú ý.

Độ phức tạp của truy hồi tính $f$ là $\Oh(k\log^2 R)$. Số cây con cần hỏi là
$\Oh(\log R)$, nên tổng độ phức tạp do phần truy hồi chi phối, tức
$\Oh(k\log^2 R)$. Nên cài đặt truy hồi bằng tìm kiếm có nhớ.

\section{Ví dụ 5: Graduated Lexicographical Ordering (ZJU 2599)}

\subsection{Đề bài}

Xét mọi số nguyên từ $1$ đến $n$. Gọi $w(x)$ là tổng các chữ số thập phân của
số nguyên $x$.

Ta sắp xếp lại các số như sau. Với hai số nguyên $a$ và $b$, nếu
$w(a)<w(b)$ thì $a$ đứng trước $b$. Nếu $w(a)=w(b)$ thì $a$ đứng trước $b$ khi
và chỉ khi biểu diễn thập phân của $a$ nhỏ hơn biểu diễn thập phân của $b$
theo thứ tự từ điển. Ví dụ:
\[
  120\grlex 4
  \quad\text{vì}\quad
  w(120)=1+2+0=3<4=w(4),
\]
\[
  555\grlex 78
  \quad\text{vì}\quad
  w(555)=15=w(78)
  \text{ và ``555'' nhỏ hơn ``78'' theo từ điển},
\]
\[
  20\grlex 200
  \quad\text{vì}\quad
  w(20)=2=w(200)
  \text{ và ``20'' nhỏ hơn ``200'' theo từ điển}.
\]

Cho $n$ và $k$, cần tìm thứ hạng của số $k$ trong dãy các số từ $1$ đến $n$
sắp theo thứ tự trên, đồng thời tìm số đứng ở vị trí thứ $k$.

Dữ liệu gồm nhiều bộ, mỗi bộ có hai số nguyên $n,k$; dòng cuối cùng là
$n=k=0$. Mỗi bộ in ra hai số nguyên: thứ hạng của $k$, và số đứng ở vị trí
thứ $k$. Ràng buộc:
\[
  1\le k\le n\le 10^{18}.
\]

\subsection{Phân tích}

Trước hết xét câu hỏi thứ nhất: cho một số $k$, tìm thứ hạng của nó. Mọi số có
tổng chữ số nhỏ hơn $\operatorname{sumof}(k)$ đều đứng trước $k$, nên ta cộng
số lượng các số này.
Vì vậy cần một hàm $\operatorname{count}(n,i)$, đếm trong $[1,n]$ có bao nhiêu
số có tổng chữ số bằng $i$. Cách cài đặt tương tự các ví dụ trước nên không
nhắc lại chi tiết.

Sau đó, ta cần đếm trong $[1,n]$ các số có tổng chữ số bằng
$\operatorname{sumof}(k)$ và nhỏ hơn $k$ theo thứ tự từ điển.

Với hai số có cùng số chữ số, thứ tự từ điển chính là thứ tự theo giá trị.
Nhưng vì số không được có chữ số $0$ ở đầu, so sánh các số khác độ dài sẽ rất
phiền. Để tránh điều này, ta chỉ so sánh các số có cùng độ dài: lần lượt xét
$k$, $k\cdot 10$, $k\cdot 100,\ldots$ và $k/10$, $k/100,\ldots$; tức là liệt
kê mọi độ dài có thể của $k$ trong phạm vi $1$ đến $n$ bằng cách thêm chữ số
$0$ hoặc bỏ chữ số ở cuối.

Với mỗi $k'$ ở một độ dài cố định, số lượng số có cùng độ dài, nhỏ hơn $k'$ và
có tổng chữ số bằng $\operatorname{sumof}(k)$ là
\[
  \operatorname{count}(k',\operatorname{sumof}(k))
  -f[\operatorname{lengthof}(k')-1,\operatorname{sumof}(k)].
\]
Ở đây $\operatorname{sumof}(k)$ là tổng chữ số của $k$,
$\operatorname{lengthof}(k')$ là độ dài thập phân của $k'$, và $f[i,j]$ là dữ
liệu tiền xử lý khi tính $\operatorname{count}(n,i)$: số lượng số trong đoạn
$[0,10^i-1]$ có tổng chữ số bằng $j$. Như vậy ta giải được câu hỏi thứ nhất.

Tiếp theo xét câu hỏi thứ hai. Lần lượt trừ khỏi $k$ các giá trị
\[
  \operatorname{count}(n,1),\operatorname{count}(n,2),\ldots
\]
cho đến khi $k$ sắp nhỏ hơn $0$; khi đó xác định được tổng chữ số $s$ của đáp
án. Việc còn lại là tìm số có tổng chữ số bằng $s$ và có thứ tự từ điển thứ
$k$ trong $[1,n]$.

Đáng tiếc là vì thứ tự từ điển khác thứ tự số học, ta không thể dùng câu hỏi
thứ nhất để nhị phân. Có thể dùng một phương pháp xác định từng chữ số khá
trực tiếp: xác định đáp án từ trái sang phải. Vì chưa biết độ dài đáp án, ta
trước hết xác định chữ số đầu tiên bằng cách liệt kê chữ số đầu và dùng thuật
toán của câu hỏi thứ nhất để kiểm tra số lượng có vượt quá $k$ hay không. Sau
đó mỗi lần thêm một chữ số vào bên phải đáp án hiện có. Nếu sau khi thêm, số
lượng số hợp lệ đúng bằng $k$ thì trả về đáp án hiện tại; nếu không thì tiếp
tục thêm cho đến khi tìm được nghiệm. Vì đáp án chắc chắn tồn tại, quá trình
thêm chữ số không vượt quá $\Oh(\log n)$ lần.

\section{Tổng kết}

Qua các ví dụ trên có thể thấy, các thao tác cơ bản trong bài toán chữ số khá
tương tự nhau và có tính mô hình rõ rệt. Trên nền các thao tác cơ bản này có
thể biến hóa thành nhiều kiểu bài khác nhau.

Tư tưởng cốt lõi khi giải là xác định từng chữ số. Cách suy nghĩ thông thường
là xét trực tiếp trên các chữ số. Tương ứng với nó là cách nhìn bằng cây, xem
các số như những lá của một cây đầy đủ. Cách thứ nhất trừu tượng hơn, còn cách
thứ hai cụ thể hơn. Vì vậy khi đề bài hơi rườm rà, dùng cách nhìn bằng cây
thường có thể biến trừu tượng thành cụ thể và làm mạch suy nghĩ rõ ràng hơn.

Do các thao tác cơ bản có độ phức tạp cấp $\Oh(\log n)$, khi xử lý một số bài
phức tạp hơn ta có thể hy sinh một phần thời gian, dùng nhị phân hoặc liệt kê
cho một vài bài toán con để giảm độ phức tạp suy nghĩ và lập trình.

Ở đây tôi chọn vài ví dụ khá tiêu biểu để giới thiệu các bài toán thống kê
theo chữ số, hy vọng có thể gợi ý cho người đọc. Cảm nhận và kỹ xảo sâu hơn
vẫn cần được tích lũy dần trong quá trình làm bài.

\section*{Lời cảm ơn}

Xin cảm ơn thầy Tạng Phương Thanh của Trường Trung học số 2 Thanh Đảo đã hướng
dẫn tôi. Xin cảm ơn huấn luyện viên Đường Văn Bân đã cung cấp ví dụ 4, bạn
Đổng Hoa Tinh của Trường Trung học số 1 Thiệu Hưng đã cung cấp ví dụ 2, và
tất cả thầy cô, bạn học đã quan tâm, ủng hộ tôi.

\appendix
\section{Đề gốc của ví dụ 1: Amount of degrees}

Hãy viết chương trình xác định số lượng số nguyên nằm trong tập $[X,Y]$ và là
tổng của đúng $K$ lũy thừa nguyên khác nhau của $B$.

Ví dụ, với $X=15$, $Y=20$, $K=2$, $B=2$, có ba số là tổng của đúng hai lũy
thừa nguyên khác nhau của $2$:
\[
  17=2^4+2^0,\qquad
  18=2^4+2^1,\qquad
  20=2^4+2^2.
\]

\textbf{Input.} Dòng đầu chứa hai số nguyên $X$ và $Y$, cách nhau bởi một dấu
cách:
\[
  1\le X\le Y\le 2^{31}-1.
\]
Hai dòng tiếp theo chứa $K$ và $B$:
\[
  1\le K\le 20,\qquad 2\le B\le 10.
\]

\textbf{Output.} In ra một số nguyên duy nhất: số lượng số nguyên nằm giữa
$X$ và $Y$ và là tổng của đúng $K$ lũy thừa nguyên khác nhau của $B$.

\textbf{Ví dụ.}
\begin{verbatim}
Input:
15 20
2
2

Output:
3
\end{verbatim}

\section{Đề gốc của ví dụ 2: Sorted bit sequence}

Xét biểu diễn $32$ bit của mọi số nguyên $i$ từ $m$ đến $n$, kể cả hai đầu
($m\le i\le n$, $mn\ge 0$,
$-2^{31}\le m\le n\le 2^{31}-1$). Số âm được biểu diễn bằng mã bù hai $32$
bit: dãy $32$ bit của nó cộng nhị phân với biểu diễn $32$ bit của số dương
tương ứng đúng bằng $2^{32}$, tức
\[
  1\,0000\,0000\,0000\,0000\,0000\,0000\,0000\,0000_2.
\]

Chẳng hạn, biểu diễn $32$ bit của $6$ là
\begin{verbatim}
0000 0000 0000 0000 0000 0000 0000 0110
\end{verbatim}
và biểu diễn $32$ bit của $-6$ là
\begin{verbatim}
1111 1111 1111 1111 1111 1111 1111 1010
\end{verbatim}
vì tổng nhị phân của hai dãy này bằng $2^{32}$.

Sắp xếp các biểu diễn $32$ bit theo thứ tự tăng dần của số bit $1$. Nếu hai
biểu diễn có cùng số bit $1$, sắp xếp chúng theo thứ tự từ điển. Với
$m=0,n=5$, kết quả là bảng đã nêu trong phần ví dụ 2; với $m=-5,n=-2$, kết quả
cũng như bảng trong phần đó.

Cho $m,n,k$ với
\[
  1\le k\le \min\{n-m+1,2147473547\},
\]
hãy tìm số tương ứng với biểu diễn thứ $k$ trong dãy đã sắp xếp.

\textbf{Input.} Dữ liệu gồm nhiều bộ test. Dòng đầu chứa số bộ test, là một số
nguyên dương không lớn hơn $1000$. Mỗi dòng sau mô tả một bộ test bằng ba số
nguyên $m,n,k$ cách nhau bởi dấu cách.

\textbf{Output.} Với mỗi bộ test, in ra trên một dòng số thứ $k$ trong dãy đã
sắp xếp.

\textbf{Ví dụ.}
\begin{verbatim}
Sample input:
2
0 5 3
-5 -2 2

Sample output:
2
-5
\end{verbatim}

\section{Đề gốc của ví dụ 3: Sequence}

Cho dãy gồm mọi số nhị phân $K$ bit:
\[
  0,1,\ldots,2^K-1.
\]
Cần chia toàn bộ dãy thành $M$ khúc. Mỗi khúc phải là một dãy con liên tiếp của
dãy ban đầu. Gọi $S_i$ ($1\le i\le M$) là tổng số bit $1$ trong mọi số của
khúc thứ $i$ khi viết ở dạng nhị phân, và gọi $S$ là giá trị lớn nhất trong
các $S_i$. Mục tiêu là tối thiểu hóa $S$.

\textbf{Input.} Dòng đầu chứa hai số $K$ và $M$:
\[
  1\le K\le 100,\qquad 1\le M\le 100,\qquad M\le 2^K.
\]

\textbf{Output.} In ra một dòng chứa giá trị nhỏ nhất có thể đạt được của
$S$, không có số $0$ thừa ở đầu. Kết quả không bảo đảm vừa trong số nguyên
$64$ bit.

\textbf{Ví dụ.}
\begin{verbatim}
Input:
3 4

Output:
4
\end{verbatim}

\section{Đề gốc của ví dụ 4: Tickets}

Nghề bán vé khá đơn điệu, vì công việc chỉ là bán vé cho hành khách. Vì vậy,
một người bán vé quyết định thay đổi công việc của mình: giờ đây anh ta bán
nhiều vé cùng lúc cho mỗi hành khách. Cụ thể, anh ta lần lượt đưa vé cho một
hành khách cho đến khi tổng các chữ số trên mọi vé đã đưa không nhỏ hơn một số
nguyên $k$. Sau đó quá trình lặp lại với hành khách tiếp theo.

Ban đầu người bán vé có một dải vé được đánh số liên tiếp từ $l$ đến $r$, kể
cả hai đầu. Cách phát vé này khiến hành khách vui vì được nhận nhiều vé dù chỉ
trả tiền cho một vé, nhưng có một bất lợi: vì mỗi hành khách nhận nhiều vé,
người bán vé có thể không phục vụ được tất cả hành khách. Hãy tính người bán
vé có thể phục vụ bao nhiêu hành khách.

\textbf{Input.} Tệp vào chứa ba số nguyên $l,r,k$:
\[
  1\le l\le r\le 10^{18},\qquad 1\le k\le 1000.
\]

\textbf{Output.} In ra đúng một số: đáp án của bài toán.

\textbf{Ví dụ.}
\begin{verbatim}
Sample input:
40 218 57

Sample output:
29
\end{verbatim}

\section{Đề gốc của ví dụ 5: Graduated Lexicographical Ordering}

Xét các số nguyên từ $1$ đến $n$. Gọi tổng chữ số của một số nguyên là trọng
số của nó, ký hiệu trọng số của $x$ là $w(x)$.

Ta sắp xếp các số theo thứ tự gọi là \term{graduated lexicographical order},
viết tắt là thứ tự grlex. Với hai số nguyên $a$ và $b$, nếu $w(a)<w(b)$ thì
$a$ đứng trước $b$ trong thứ tự grlex. Nếu $w(a)=w(b)$, thì $a$ đứng trước $b$
trong thứ tự grlex khi và chỉ khi biểu diễn thập phân của $a$ nhỏ hơn biểu
diễn thập phân của $b$ theo thứ tự từ điển.

Ví dụ:
\[
  120\grlex 4
  \quad\text{vì}\quad
  w(120)=1+2+0=3<4=w(4),
\]
\[
  555\grlex 78
  \quad\text{vì}\quad
  w(555)=15=w(78)
  \text{ và ``555'' nhỏ hơn ``78'' theo từ điển},
\]
\[
  20\grlex 200
  \quad\text{vì}\quad
  w(20)=2=w(200)
  \text{ và ``20'' nhỏ hơn ``200'' theo từ điển}.
\]

Cho $n$ và một số nguyên $k$, hãy tìm vị trí của số $k$ trong thứ tự grlex của
các số nguyên từ $1$ đến $n$, đồng thời tìm số nguyên đứng ở vị trí thứ $k$
trong thứ tự này.

\textbf{Input.} Dữ liệu gồm nhiều dòng, mỗi dòng chứa hai số nguyên $n$ và
$k$:
\[
  1\le k\le n\le 10^{18}.
\]
Dòng $n=k=0$ kết thúc dữ liệu.

\textbf{Output.} Với mỗi dòng vào, in ra một dòng. Trước hết in vị trí của số
$k$ trong thứ tự grlex của các số nguyên từ $1$ đến $n$, sau đó in số nguyên
đứng ở vị trí thứ $k$ trong thứ tự này.

\textbf{Ví dụ.}
\begin{verbatim}
Sample Input:
20 10
0 0

Sample Output:
2 14
\end{verbatim}

\end{document}
